% Release candidate v1.0.0; see the W04 preparation record for artifact checks.
% G06/W03 is a separate frozen snapshot; see README.md for the evidence ledger.
\documentclass[11pt,a4paper]{article}
\usepackage{amsmath,amssymb,amsthm,array}
\usepackage[margin=27mm]{geometry}
\usepackage[T1]{fontenc}
\usepackage{lmodern,microtype,fancyhdr,xurl}
\usepackage[hidelinks]{hyperref}
\hypersetup{pdftitle={Coverage in Heraclitus walks: interval filling and residue connectivity},pdfauthor={Roger Gujord}}
\pagestyle{fancy}\fancyhf{}
\fancyhead[L]{\small Coverage in Heraclitus walks}
\fancyhead[R]{\small Preprint v1.0.0}
\fancyfoot[C]{\thepage}\setlength{\headheight}{14pt}
\setlength{\parskip}{4pt}
\setlength{\parindent}{0pt}
\widowpenalty=10000
\clubpenalty=10000
\displaywidowpenalty=10000
\emergencystretch=1em
\newtheorem{maintheorem}{Theorem}
\renewcommand{\themaintheorem}{\Alph{maintheorem}}
\newtheorem{theorem}{Theorem}
\newtheorem{lemma}[theorem]{Lemma}
\newtheorem{corollary}[theorem]{Corollary}
\newtheorem{proposition}[theorem]{Proposition}
\newcommand{\Z}{\mathbb Z}\newcommand{\N}{\mathbb N}
\newcommand{\floor}[1]{\left\lfloor #1\right\rfloor}
\newcommand{\ceil}[1]{\left\lceil #1\right\rceil}
\newcommand{\code}[1]{\texttt{#1}}
\title{Coverage in Heraclitus walks\\[5pt]\large Interval filling and residue connectivity}
\author{Roger Gujord\\\small ORCID: \href{https://orcid.org/0009-0006-5227-980X}{0009-0006-5227-980X}}
\date{5 November 2026\\\small Version 1.0.0\\[4pt]\footnotesize
DOI: \href{https://doi.org/10.5281/zenodo.23973601}{10.5281/zenodo.23973601}}
\begin{document}
\maketitle
\begin{abstract}
A Heraclitus walk starts at zero and repeatedly chooses the unused integer
closest to zero at an allowed distance, preferring the positive value in a
tie. This paper proves two sufficient conditions for visiting every integer
exactly once. Every unbounded distance set containing 1 and 2 has coverage,
with exact complete-interval checkpoints at every allowed distance at least
2. Coverage also holds when the distances contain complete tails in a
nonempty collection of residue classes whose representatives, together with
the modulus, have greatest common divisor 1. The latter condition requires
no prescribed short distances. The proofs use interval filling and a
finite-priority argument that propagates omitted targets through residue
classes. For a full odd tail above a threshold $M\ge1$, the first visit
to $z$ occurs by index $4M+6|z|+9$. A covered family outside both
criteria shows that they are not exhaustive. These results answer Sloane's general coverage question for two classes of
distance sets and unify previously reported special cases. The main results
are formalized in Lean 4.35.0 using the rule stated here. The verification
scope and prior-work attribution are given below.
\end{abstract}

\section{Introduction and main results}
Throughout, $\N=\{0,1,2,\ldots\}$. The set $D\subseteq\N$ is fixed for the entire
walk. It specifies allowed \emph{distances}, not a prescribed chronological
list of jumps. Put $a(0)=0$ and $U_n=\{a(0),\ldots,a(n)\}$. At time $n$,
the next point is the least admissible unused destination in the order
\[
0,1,-1,2,-2,3,-3,\ldots,
\qquad v\text{ admissible}\iff |v-a(n)|\in D.
\]
Equivalently, minimize $(|v|,\mathbf1_{v<0})$ lexicographically over admissible
$v\notin U_n$. Allowing $0\in D$ has no effect: the current point is used.

The name \emph{Heraclitus walk} follows Sloane's terminology for this
greedy construction.
Sloane asks whether an infinite set of positive allowed distances with
greatest common divisor 1 must have a Heraclitus transform containing every
integer \cite{sloane2025,sloane2026}. Both conditions are necessary. A common
divisor confines the walk to its multiples. A finite distance set bounds
all jumps, while an injective walk eventually leaves every finite interval;
it can then no longer cross between the two unbounded sides. The question
is whether the two conditions are sufficient.

The two theorems below prove sufficiency for two general classes. The first permits arbitrarily
sparse large distances but requires 1 and 2. The second requires complete
residue tails but permits all small distances to be absent. A full tail in
a residue class contains every sufficiently large positive integer in that
class. A checkpoint specifies the visited set and current point at a given
time. To state the exact checkpoints of the first theorem, define
\begin{align}
 I_n&=[-\floor{n/2},\ceil{n/2}]\cap\Z,\label{eq:box}\\
 E_n&=\{\ceil{n/2}-1,\ 1-\floor{n/2},\ \ceil{n/2},\ -\floor{n/2}\}.
 \label{eq:ends}
\end{align}
Some elements of $E_n$ coincide for small $n$; it is a set, not a claim
that four distinct states occur in every walk.

\begin{maintheorem}[The 1--2 coverage theorem]\label{thm:main}
Suppose $1,2\in D$ and $D$ is unbounded: for every $N$ there is $d\in D$
with $d>N$. A unique infinite walk following the stated rule exists. Every such
walk satisfies, for every $d\in D$ with $d\ge2$,
\[
 U_d=I_d,\qquad a(d)\in E_d.
\]
Moreover, every $z\in\Z$ is attained at exactly one index.
\end{maintheorem}

\begin{maintheorem}[Residue-tail coverage]\label{thm:residue-tail}
Let $q\ge1$ and $M\ge0$ be integers. Let $S$ be a nonempty subset of
$\{0,\ldots,q-1\}$. Suppose the greatest common divisor of $q$ and all
elements of $S$ is 1, written
\[
 \gcd(q,\{s:s\in S\})=1.
\]
Suppose that every positive $d\ge M$ with $d\bmod q\in S$ belongs to $D$.
Other allowed distances are arbitrary. The walk starting at zero and following the stated rule
exists uniquely and visits every integer exactly once.
\end{maintheorem}

For example, Theorem~\ref{thm:main} applies to every unbounded set after
adding 1 and 2. Theorem~\ref{thm:residue-tail} applies to the exact set
$\{5,8,11,14,\ldots\}$, with modulus 3 and residue 2.
Theorems~\ref{thm:main} and~\ref{thm:residue-tail} concern a fixed set of
distances; adding distances can change the walk.
Neither theorem asserts that a prescribed sequence of jump lengths is used.

The first proof fills intervals and tracks their endpoints. The second
uses the fact that repeatedly passing over a missing target requires
choosing a different point each time from a finite set of destinations that
take priority over it. These
mechanisms give coverage under different hypotheses. A separate counting
argument supplies the odd-tail first-visit bound in
Proposition~\ref{prop:odd-tail}.

The two criteria are not exhaustive. The covered reference set
$\{1,3,6,10,14,\ldots\}$ in Appendix~\ref{sec:triple} satisfies neither:
it omits 2, and all its sufficiently large distances are even. The broader
coverage question therefore remains beyond the scope of these criteria.

The accompanying GitHub repository contains the Lean proofs, a guide linking
the paper's statements to their formal counterparts, and instructions for
reproducing the verification. The versioned repository link and archived
source package \cite{package} are described in Section~\ref{sec:formal}.

\section{Related work}\label{sec:context}
The following OEIS reports were checked on 5 November 2026.
Several special cases were already reported. Sigrist gives a block
description for powers of two \cite{powers}. Caveney announces a Fibonacci
coverage proof and a six-block outline \cite{fibonacci}, and a
squares-plus-2 proof with a four-block outline \cite{squares2}.
His companion entry explicitly explains why adding 2 simplifies the
problem \cite{adding2}. The factorial project of trureturing/The Omega
Institute supplies a formal source and an interval-filling argument
\cite{factorial,factorialcode}. The inward descent and return-filling argument used here follows that
factorial work.
Neither the filling idea nor the already reported named coverage results
are presented as first discoveries of this project.

Theorem~\ref{thm:main} treats this filling mechanism for arbitrary unbounded
distance sets containing 1 and 2, with exact complete-interval checkpoints
and endpoint constraints at every allowed distance at least 2.
Theorem~\ref{thm:residue-tail} gives a complementary criterion using complete
residue tails with joint gcd 1, without requiring any particular short distance.

The available Caveney outlines are credited here; his complete SeqFan
correspondence has not been retrieved. The factorial source was inspected
at the pinned commit in the bibliography, without an external build or
axiom audit. Priority of the arbitrary-set statement, the stronger
all-allowed-distance checkpoints and the residue-tail criterion must be
assessed separately. The source comparison does not establish first
discovery or exclude unpublished work.

\section{Existence, uniqueness of choices, and exact indices}
For a finite used set $U$ and current point $c$, choose $B$ with
$|u|\le B$ for all $u\in U$. Unboundedness gives $d\in D$ with
$d>B+|c|$. Then $c+d>B$ is unused and reachable. Thus there is always an
admissible unused destination. The displayed order on $\Z$ is a well-order
of type $\omega$, so the next destination exists and is unique.
Recursive choice on finite states gives a unique infinite walk. Any finite
sequence of greedy moves from a state reached by that walk agrees with its
continuation, since the next choice is unique. For arbitrary $D$, a next
choice is unique whenever an admissible unused destination exists;
unboundedness ensures existence at every step.

No step repeats a point, hence $|U_n|=n+1$. In particular, $I_m$ has $m+1$
points. If a finite greedy continuation from a state reached by the walk
has visited exactly $I_m$, its ending time is $m$. This counting argument
determines the time whenever a local transition fills such an interval;
an endpoint position alone
does not determine the time or the visited set.

\begin{lemma}[Finite priority for an omitted target]\label{lem:finite-priority}
Let $a$ be an infinite walk following the stated rule with distance set $D$.
If $h$ is never visited, then there exists $N$ such that
\[
 |h-a(n)|\notin D\qquad\hbox{for every }n\ge N.
\]
\end{lemma}
\begin{proof}
Whenever $h$ is reachable from the current point, it is an admissible unused
competitor. Since $h$ is never selected, the chosen successor precedes it
in the priority order and lies in $[-|h|,|h|]$. These successors are
distinct. Hence only finitely many indices can have $h$ reachable, which
gives the assertion. The comparison allows equal absolute values: a
positive successor can precede a negative target of the same magnitude.
\end{proof}

The same successor comparison holds while a target is absent from a finite
prefix. Counting the distinct small successors, rather than only proving
eventual avoidance, is the ingredient used for the time bound in
Section~\ref{sec:odd-tail}.

\section{Filling one side}\label{sec:filling}
The following version states explicitly what is needed about the exterior.

\begin{lemma}[Interior filling]\label{lem:fill}
Let $1\le L<R$, $\varepsilon\in\{-1,1\}$, and $1,2\in D$.
Suppose the current point is $\varepsilon R$, and the unused points with
absolute value strictly less than $R$ are precisely
\[
 \{\varepsilon t:L\le t<R\}.
\]
Then finitely many greedy moves fill exactly these holes and end at
$\varepsilon(R-1)$. All these moves have absolute destination value less
than $R$. Which points outside that strict interior have been visited does not affect
the conclusion.
\end{lemma}

\begin{proof}
Work with magnitudes on the active side; a magnitude is \emph{marked} when
its point on that side has already been visited. During inward descent, let the
current magnitude be $r$. All magnitudes in $[L,r)$ are unvisited, while
$r$ and $R$ are marked. Maintain that no two marked magnitudes $t,t+1$
have $L<t<R$; the possible pair $L,L+1$ is exempt. Initially only $R$ is
marked, so this invariant holds.

If $r>L$, distance 1 reaches the unused magnitude $r-1$. The greedy
destination must therefore have magnitude $s$ with $L\le s<r$, on the
active side. If $r\ge L+2$, distance 2 reaches $r-2$, so $s\le r-2$.
If $r=L+1$, necessarily $s=L$. A new mark $s>L$ has both neighbors
unmarked, because they lay below the previous current. A new mark at $L$
can create only the exempt bottom adjacency. Thus the invariant persists,
and strict decrease reaches $r=L$.

For the return, maintain that all active magnitudes through $r$ are used,
that $R$ is marked, and that there are no adjacent marked magnitudes
$t,t+1$ with $r<t<R$. If $r=R-1$, filling is complete. Otherwise, if
$r+1$ is unused, it is the least-magnitude unused point and distance 1
reaches it. If $r+1$ is used, the sparsity invariant makes $r+2$ unused;
since $R$ is used, necessarily $r+2<R$. Distance 2 reaches this next
least-magnitude hole. Both choices preserve the return invariant and
increase $r$ while keeping $r\le R-1$. The return ends at $R-1$.

All selected destinations have magnitude below $R$, where the opposite
side is already occupied. Therefore no exterior point or opposite-sign
tie can change the argument. No reflection symmetry of the tie rule is
being assumed.
\end{proof}

When $L=R$, the hole interval is empty: a crossing that has just visited
$\varepsilon R$ already completes this side. There is no filling move and
the endpoint remains $\varepsilon R$. Keeping this case is the reason for
the boundary states in $E_n$.

\section{Transitions between consecutive allowed distances}
Let $q<r$ belong to $D$, with $q\ge2$ and no member of $D$ strictly between
them. Suppose the current used set is $I_q$ and the current point lies in
$E_q$. The following argument gives a finite greedy continuation to $I_r$
with endpoint in $E_r$.

\subsection{Run, crossing, and filling when \texorpdfstring{$r\ge q+2$}{r >= q+2}}
Put $B=\floor{r/2}$ and $T=\ceil{r/2}$, so $B\le T\le B+1$ and
$B+T=r$. The proof uses two elementary blocks of moves.

For a \emph{positive block}, suppose the state is current $x$ with used
set $[-b,x]\cap\Z$, where
\[
 q<2b+2,\quad b+1\le x\le T,\quad b+1\le B.
\]
If $x<T$, the point $x+1$ is reachable by 1. Any unused destination of
smaller magnitude must be $-t$, with $b+1\le t\le x$, and its distance
from $x$ satisfies
\[
 q<2b+2\le x+t\le2x\le2T-2<r.
\]
Such a move is forbidden. A negative destination tied in magnitude loses
to the positive candidate, whether or not it is reachable. Thus unit
moves reach $+T$. From there, $-B$ is reachable by distance $r$. Any
smaller unused competitor is $-t$, $b+1\le t<B$, and
\[
 q<2b+2\le T+b+1\le T+t<T+B=r.
\]
All positive points of magnitude at most $B$ are used. Hence the crossing
to $-B$ is greedy. The remaining holes are the negative magnitudes
$b+1,\ldots,B-1$. Lemma~\ref{lem:fill}, or its empty-hole case, gives
the full interval $[-B,T]\cap\Z$ and endpoint $1-B$ or $-B$.

For a \emph{negative block}, suppose the state is current $-x$ with used
set $[-x,u]\cap\Z$, where
\[
 q<2u+1,\quad u\le x\le B,\quad u+1\le T.
\]
For $x<B$, an unused competitor beating \emph{or tying} $-(x+1)$ must
be positive $t$, with $u+1\le t\le x+1$. Its distance satisfies
\[
 q<2u+1\le x+t\le2x+1\le2B-1<r.
\]
It is forbidden, including the positive tie that would otherwise win.
Unit moves therefore reach $-B$. The $r$-crossing reaches $+T$. A smaller
unused positive competitor $t$ has $u+1\le t<T$, giving
$q<2u+1\le B+t<r$. All smaller unused negative competitors are excluded
because those points have been visited; when $T=B+1$, the possible point
$-T$ loses the tie.
Filling positive magnitudes $u+1,\ldots,T-1$ now ends at $T-1$, or at
$T$ when the crossing already filled the last hole.

These blocks start from all four checkpoint states as follows. Write
$q=2b$ or $q=2b+1$.
\begin{center}
\renewcommand{\arraystretch}{1.25}
\begin{tabular}{c|c|c|c}
$q$ & current point & preparation & block\\\hline
$2b$ & $b-1,\ 1-b,\ b$ & go to $b+1$ & positive\\
$2b$ & $-b$ & none & negative, $u=x=b$\\
$2b+1$ & $b,\ 1-b,\ -b$ & go to $-(b+1)$ & negative\\
$2b+1$ & $b+1$ & none & positive, $x=b+1$
\end{tabular}
\end{center}
In the even first row, the required distances are respectively $2,q,1$;
the positive destination wins the least-unused-magnitude tie. In the odd
third row, the required distances are $q,2,1$; $+(b+1)$ is already used,
so $-(b+1)$ is the least unused point. The inequalities for the indicated
blocks follow directly from $r\ge q+2$. In each case their output is
$I_r$ with endpoint in $E_r$. This also covers $q=2$ and $q=3$.

\subsection{Adjacent allowed distances}
If $r=q+1$, exactly one new point is needed. For even $q=2b$ it is
$b+1$; from the ordered endpoints $b-1,1-b,b,-b$, the distances are
$2,q,1,q+1$. All are allowed, and the positive destination is the preferred
least-magnitude unused point. For odd $q=2b+1$, the new point is $-(b+1)$;
from endpoints $b,1-b,b+1,-b$, the distances are $q,2,q+1,1$.
The positive point $b+1$ is already used, so again the move is greedy.
The new endpoint is a boundary point of $I_{q+1}$.

This proves the following local statement, with no restrictions on the
other distances below $q$ or above $r$.
\begin{lemma}[Checkpoint transition]\label{lem:transition}
If $1,2,q,r\in D$, $2\le q<r$, and $D\cap(q,r)=\varnothing$, then
every state with used set $I_q$ and current in $E_q$ has a finite greedy
continuation to used set $I_r$ and current in $E_r$.
\end{lemma}

\section{Assembly and coverage}
The first two moves are $+1$ and $-1$, reached using 1 and 2. Thus
$U_2=I_2=[-1,1]\cap\Z$ and $a(2)=-1\in E_2$.
Let $s_0=2$ and let $s_{j+1}$ be the least member of $D$ greater than $s_j$.
Unboundedness ensures this exists. The sequence is a construction in the
proof; it requires no algorithm for deciding membership in $D$. It satisfies
$s_j\ge j+2$ and $D\cap(s_j,s_{j+1})=\varnothing$.

Inductively apply Lemma~\ref{lem:transition}. Its finite path agrees with
the walk starting at zero. Counting the visited points puts its last state at
index $s_{j+1}$, so $U_{s_j}=I_{s_j}$ and $a(s_j)\in E_{s_j}$ for all $j$.
Every allowed $d\ge2$ occurs in this enumeration: take the first $s_j\ge d$.
If $j=0$, then $d=2$. Otherwise, $s_{j-1}<d\le s_j$, and strict inequality
$d<s_j$ would contradict the least-successor definition. This proves all
the checkpoints in Theorem~\ref{thm:main}.

For $z\in\Z$, use the checkpoint at $s_{2|z|}$. Its index is at least
$2|z|+2$, so its interval contains $z$. No point is revisited, so each occurrence is unique.
Together with the existence construction, this proves the theorem.

\begin{corollary}[Completion by two short distances]\label{cor:completion}
For every unbounded $S\subseteq\N$, the fixed distance set
$S\cup\{1,2\}$ has an integer-covering greedy walk with the checkpoints
in Theorem~\ref{thm:main}.
\end{corollary}
Adding a missing distance changes the rule and generally changes the walk.
This corollary does not transfer coverage back to the unmodified set.

\section{Coverage from full residue tails}\label{sec:residue-tail}
The complete-tail assumption in Theorem~\ref{thm:residue-tail} turns
finite priority into a constraint on entire residue classes. Restricting
allowed distances to those classes, or having infinitely many distances in
them, would not supply this implication.

\begin{proof}[Proof of Theorem~\ref{thm:residue-tail}]
A residue in $S$ supplies arbitrarily large allowed distances, giving
existence and uniqueness by the finite-history argument above. Let $V$
be the set of visited integers and suppose $h\notin V$.
Lemma~\ref{lem:finite-priority} implies that
\[
 V\cap(h+D)\quad\hbox{and}\quad V\cap(h-D)
\]
are finite.

Fix $s\in S$. Every sufficiently large positive integer congruent to
$h+s$ modulo $q$ lies in $h+D$, so only finitely many visited positive points lie in that residue class. Likewise every sufficiently negative integer congruent to
$h-s$ lies in $h-D$. Each of these two parts of a residue class therefore contains unvisited points.

Let $W$ be the nonempty set of residues containing an omitted integer.
The preceding argument gives $W+s\subseteq W$ for every $s\in S$.
Translation is a bijection of the finite residue space, so $W+s$ and $W$
contain the same number of elements; hence $W+s=W$. Translation by $q$ is the
identity modulo $q$. The integer shifts preserving $W$ are
closed under differences and integer multiples. Euclid's algorithm
therefore shows that translation by $\gcd(q,S)=1$ preserves $W$.
Since $W$ is nonempty, it contains every residue.

Choose one $s\in S$. For each residue $b$, choose an omitted target of
residue $b-s$ and one of residue $b+s$. The first target bounds the visited
positive points of residue $b$; the second bounds the visited negative
points there. Hence each residue contains only finitely many visited
integers. There are $q$ residues, so $V$ would be finite, contradicting the
infinite walk with no repeated values. Thus every integer occurs exactly once.
\end{proof}

For a single residue $r$, the condition is simply $\gcd(q,r)=1$.
For example, the exact distance set $\{5,8,11,14,\ldots\}$ is covered,
although it contains none of 1,2,3. Several residues can generate jointly
even when none is individually coprime to the modulus: the distances
$d\ge8$ with $d\bmod6\in\{2,3\}$ are another covered example.
For modulus one, the hypothesis simply requires all sufficiently large distances.

The theorem permits arbitrary additional distances and finite deletions
from the required tails, the latter by increasing $M$. It does not apply
to sparse subsets merely because their residues satisfy the gcd condition.
Its formal statement uses a finite list of canonical representatives;
duplicates have no effect on the argument. Nonemptiness is necessary even
when $q=1$, since an empty list would impose no distance condition.

The residue-tail proof gives no explicit first-visit bound. The odd-tail
bound below therefore needs a separate proof.

The two main proofs use the same priority order in different ways.
With distances 1 and 2, local descent and return fill an interval and
produce exact checkpoints, however sparse the larger distances are.
With full residue tails, a permanently omitted target would force too
many distinct successors into a finite priority interval. Residue
connectivity propagates that obstruction until the entire walk would be
bounded. This second mechanism permits every short distance to be absent,
but depends on the complete tails.

\section{A quantitative bound for odd tails}\label{sec:odd-tail}
The odd-tail case of Theorem~\ref{thm:residue-tail} admits a separate
explicit time bound. Its proof is included because the qualitative residue
argument does not provide this constant.

\begin{proposition}[Cofinite odd tail and first-visit bound]\label{prop:odd-tail}
Let $M\ge1$ be an integer. Suppose every odd integer $d\ge M$ belongs
to $D$, with all other positive distances arbitrary. The walk defined above
exists, is unique, and visits every integer exactly once. Its first
visit to $z$ satisfies
\[
 \tau(z)\le 4M+6|z|+9.
\]
None of distances 1,2,3 is required. The bound is not claimed to be optimal.
\end{proposition}
\begin{proof}
The allowed odd tail is unbounded, so existence follows from the earlier
finite-state argument. The fixed priority order makes the walk unique, and
the rule forbids repeated visits.

Suppose $y$ is absent from the prefix through index $N$. Among the $N$
current points $a(0),\ldots,a(N-1)$, at most $M+2|y|+1$ have parity
opposite to $y$. Indeed, at most $M$ such distinct points lie at distance
less than $M$ from $y$. For each remaining point the odd distance to $y$
is allowed, and $y$ is unused. Greedy priority forces the successor into
$[-|y|,|y|]$. These successors are distinct, so there are at most
$2|y|+1$ of them. This proves the count without counting the endpoint
$a(N)$ as a current point.

Now suppose $z$ is absent and put $R=|z|$. Including the endpoint, the
whole prefix contains at most $M+2R+2$ points of parity opposite to $z$.
Set $B=M+2R+3$. The interval $[-B,B]$ contains at least $B$ integers of
that parity, so it contains a missing point $w$ of opposite parity to $z$,
with $|w|\le B$. Applying the count to this second missing target bounds
the current points of the same parity as $z$ by $M+2B+1$.
The two parity classes partition the $N$ current points, giving
\[
 N\le (M+2R+1)+(M+2B+1)=4M+6R+8.
\]
Thus omission at $N=4M+6R+9$ is impossible. This proves the first-visit
bound and coverage. The rule forbids repeated visits, so each integer occurs
exactly once.
\end{proof}

For the family $D_k$ in Appendix~\ref{sec:parity}, take $M=10k+3$.
Every point of $[-5k-2,5k+2]$ is
visited by $70k+33$. Consequently all holes in its state at
time $5k+5$ are repaired within a further $65k+28$ moves. Coverage itself
uses only Proposition~\ref{prop:odd-tail}; identification of those
particular holes also uses Proposition~\ref{prop:many-holes}.
This does not require repair before an excursion.

The proposition includes the rule with exactly the odd distances at least 5,
which omits all three short distances 1,2,3. Arbitrary additional distances
can alter its trajectory without invalidating the missing-target comparison.
The original prime rule omits arbitrarily large odd composites and therefore
fails the odd-tail hypothesis. Sparse odd sets, the base-seven rule and
prime-gap laws are not covered by this argument.

\section{Applications and scope}
The first theorem includes powers of two, positive Fibonacci values,
factorials, and squares together with 2. Their previously reported coverage
results and antecedent mechanisms are credited in Section~\ref{sec:context}.
Adding 1 and 2 to any unbounded set supplies the formal completion
corollary, but changes the rule; coverage does not transfer back after a
required distance is removed. No restriction on the gaps between large
allowed distances is needed.

The two general criteria do not settle the original prime, triangular or
odd-base power rules. Ordinary squares also fail their hypotheses; a separate
AI-assisted proof of the square-walk conjectures was announced on 29 October
2026 and is stated to be under review in OEIS A377091 \cite{squares}.
The announced proof has not been checked as part of this work. The separate base-seven study
\cite{base7} concerns a rule outside these criteria; its claims are not
strengthened here. Coverage of a derived walk gives no prime-gap bound by
itself, since the first theorem permits arbitrary gaps in the input set.

\section{Verification and AI contribution}\label{sec:formal}
The numbered mathematical statements have corresponding Lean 4.35.0 proofs.
The coverage theorems use the stated rule starting at zero; local transition and
filling results retain their explicit state hypotheses. The formal rule
forbids repeated visits, compares all integer candidates, and prefers the
positive candidate in a tie. Infinite-walk existence is proved separately.
The finite-priority lemma is the permanent-omission statement
\code{G18.missing\_avoids}; the quantitative proof uses its finite-prefix
successor comparison. The source package \cite{package} contains a statement-to-proof guide in
\code{docs/REVIEW\_GUIDE.md}, the complete required Lean sources, and
reproduction instructions. Its versioned GitHub location is
\url{https://github.com/gujord/heraclitus-coverage/tree/v1.0.0}.

The proofs use Lean's bundled standard library. Their transitive axioms are
restricted to \code{propext}, \code{Classical.choice} and \code{Quot.sound}.
No empirical trajectory sample is a premise of the formal results.
Trust includes the Lean kernel and the pinned toolchain. A fresh combined
build from an extracted source archive compiled all 44 selected modules,
audited all 296 retained theorems, and rejected all 40 intended false-proposition
fixtures. The package selects the declarations needed for this paper and its
regression checks, preserving their statement and proof bodies while narrowing
imports. All 402 retained declaration types agree with the earlier accepted
builds. The provenance record identifies those builds separately from this
combined reproduction. Kernel checking does not by itself establish
historical originality or the fidelity of this prose to the formal statements.

AI tools contributed to mathematical exploration and generated proof drafts,
exposition and Lean formalizations under Roger Gujord's direction. The
antecedents and inspected source material are credited in
Section~\ref{sec:context}. Roger Gujord is the sole author and is responsible
for the paper. AI assistance and local machine verification do not
constitute independent human review.
Historical originality and independent human mathematical review remain open.

\section{Scope and open questions}\label{sec:questions}
The main unresolved issue is Sloane's general sufficiency question: does
every infinite distance set with gcd 1 cover all integers?
Theorems~\ref{thm:main} and~\ref{thm:residue-tail} answer the question for two classes of distance sets, while Appendix~\ref{sec:triple}
supplies further covered sets outside them. They do not give a classification.

A more specific question asks for broader coverage conditions when 1 is
allowed and 2 is absent. The appendix family settles some such rules but
does not justify replacing the 1--2 hypothesis by 1--3 in an arbitrary
unbounded distance set. For the prime distance set, coverage remains the
conjecture recorded in OEIS A377090 \cite{primes}; its sparse odd distances
do not satisfy the complete-tail premise, and 1 is missing.

For sparse distance sets such as the primes, a possible approach is to
prove that the walk repeatedly reaches configurations from which the next
missing integer is forced to be visited. Establishing that recurrence from
the walk's full history remains an open proof obligation.

\paragraph{Moving checkpoint holes.}\label{rem:moving-hole}
Exact complete intervals are sufficient for the first proof but are not
necessary for coverage. The checked example
$D=\{1,3,10\}\cup\{64+4j:j\ge0\}$ has, for every $j\ge0$,
at $n=72+4j$ and $b=36+2j$,
\[
 a(n)=-b,\qquad U_n=([-b,b-1]\cap\Z)\cup\{b+1\},\qquad a(n+1)=b.
\]
The complete interior grows without bound while a boundary hole is repeatedly
created and repaired. Appendix~\ref{sec:repair} gives the finite initial-path certificate and the
four-move induction proving these formulas and coverage. Appendix~\ref{sec:parity} gives a different limitation: there
is no uniform number of lane holes across its family, even though every
member has coverage by Proposition~\ref{prop:odd-tail}. These observations
restrict particular state descriptions; they are not counterexamples to
Sloane's question.

\appendix

\section{A coverage family without distance 2}\label{sec:triple}
This family demonstrates that Theorems~\ref{thm:main} and
\ref{thm:residue-tail} are not exhaustive. Its reference set omits 2 and
has only finitely many odd distances. Any full arithmetic tail contained
in such a set must have both even modulus and even residue, so a collection
of those tails cannot satisfy the joint-gcd condition. Optional extensions
of the reference set need not share this exclusion.

\begin{theorem}[A family without distance 2]\label{thm:triple}
Let
\[
 D=\{1,3\}\cup\{6+4j:j\ge0\}\cup T,
 \qquad T\subseteq\{4\}\cup\{7+4j:j\ge0\}.
\]
An infinite walk following the stated rule exists, and every such walk
visits each integer exactly once. For every $j\ge0$,
\[
 U_{6+4j}=[-3-2j,3+2j]\cap\Z.
\]
Moreover $a(6)=3$, and $a(10+4j)=-4-2j$ for $j\ge0$.
The subset $T$ is arbitrary, with no computability or periodicity hypothesis.
\end{theorem}

\begin{proof}
The initial points are
\[
 0,1,2,-1,-2,-3,3.
\]
From 1, the candidate $-1$ requires the absent distance 2; $+2$ wins
the next sign tie. Distance 3 then reaches $-1$, and distance 1 reaches
$-2$. From $-2$, $+3$ would require the absent distance 5, so $-3$ is
next. Distance 6 reaches $+3$. Optional 4 and larger optional distances
cannot beat these choices. Thus time 6 has used interval $[-3,3]$ and
current point $+3$.

For the induction, suppose the used interval is $[-b,b]$, $b\ge1$, and
the current is $b$ or $1-b$. Suppose $1,2b,2b+4$ are allowed, while
$2b+2,2b+3$ are absent. The next four points are
\[
 b+1,\quad b+2,\quad -(b+2),\quad -(b+1).
\]
The first point has least unused magnitude, is reachable by 1 or $2b$,
and wins the positive tie. From $b+1$, the smaller unused point $-(b+1)$
requires the forbidden $2b+2$, so the positive unit move wins. From $b+2$,
the only unused point of smaller magnitude is $-(b+1)$, requiring the
forbidden $2b+3$. Hence the allowed $2b+4$ crossing reaches $-(b+2)$.
Finally distance 1 reaches the least-magnitude unused point $-(b+1)$.
All other unused points are farther from zero. The output is exactly
$[-b-2,b+2]$, with current $-(b+1)=1-(b+2)$.

Set $b=3+2j$. The required distances are $6+4j$ and $6+4(j+1)$;
the excluded ones are $8+4j$ and $9+4j$. Every allowed optional subset
preserves these comparisons. Four moves give the stated exact indices
and endpoints by induction. The intervals exhaust $\Z$, and no point is
revisited, so each integer occurs exactly once. Unboundedness of the mandatory progression gives
existence by the finite-state argument already proved above.
\end{proof}

An \emph{auxiliary state} specifies a current point and used set; it need
not be reached by the walk starting at zero.

\begin{proposition}[A failed local replacement]\label{prop:triple-obstruction}
Let $R=\{1,3\}\cup\{6+4j:j\ge0\}$. From the auxiliary state with
current $-8$ and used set $\{-8\}\cup([0,8]\cap\Z)$, no finite greedy
continuation has used set exactly $[-8,8]\cap\Z$.
\end{proposition}
\begin{proof}
The forced path, including its initial point, is
\[
 -8,-2,-1,-4,-3,-6,-5,9.
\]
At each inward stage the displayed destination is the first admissible unused
point in the priority order. Upon reaching $-5$, the only remaining hole in
$[-8,8]$ is $-7$, whose distance is the absent 2. The next winner is $+9$
at distance 14. Every earlier state still has an interior hole, and every
later used set contains the exterior point 9. Since used points cannot
be removed, the complete interval can never be the exact used set.
\end{proof}
This obstruction concerns only the specified starting state. Theorem~\ref{thm:triple} proves coverage for the
very same rule $R$. Nor does the obstruction assert that $-7$ is never
visited after escape. It rejects only the direct 1--3 substitution into
the one-sided filling statement in Lemma~\ref{lem:fill}.

\section{Missing distance 1: reachable holes and a limited filling lemma}
\label{sec:parity}
Both propositions below are formalized in Lean, including the exact
state reached from zero, its hole count, and the conditional local filling statement.
A parity lane is the subsequence of integers of one parity. Its visited
hull consists of all same-parity integers between its visited minimum and
maximum; a lane hole is an unused point in that hull.

\begin{proposition}[No uniform lane-hole count]\label{prop:many-holes}
For every $k\ge1$, let
\[
 D_k=\{2,3\}\cup\{10k+3+2t:t\ge0\}.
\]
At time $n=5k+5$ the walk starting at zero satisfies
\[
 a(n)=5k+1,\qquad
 U_n=([-5k-2,2]\cap\Z)\cup\{5k+1\}.
\]
The two visited parity-lane hulls together have exactly
$\floor{(5k-2)/2}$ holes. Thus no fixed constant bounds the number of
lane holes over the entire class of unbounded sets consisting of 2,3
and otherwise odd distances.
\end{proposition}
\begin{proof}
First use only the distances $F=\{2,3\}$. Its first seven points are
$0,2,-1,1,-2,-4,-6$. At time $5j+1$, $j\ge1$, write $t=5j$.
The inductive state is current $-t-1$ and used set
$([-t-1,2]\cap\Z)\setminus\{-t,-t+2\}$. Its next five points are
\[
 -t+2,\ -t,\ -t-2,\ -t-4,\ -t-6.
\]
The first fills the reachable interior hole by distance 3; the other hole
is only distance 1 away. The second fills that remaining hole by distance 2.
For each of the last three moves, both inward short-step destinations are
used, and the outward distance-2 point precedes the distance-3 point.
Checking the four candidates at each move proves the displayed state
for $j+1$ and hence the induction.

Now fix $k$ and put $t=5k$, $M=2t+3$. Through time $5k+1$, all current points
and chosen destinations in this path have magnitude at most $t+1$.
Any destination using an additional distance $d\ge M$ has magnitude
at least $d-|c|\ge t+2$, so it cannot change these choices.
The next three destinations using short steps are $-t+2,-t,-t-2$; the corresponding
large-step lower bounds are $t+2,t+5,t+3$, again strictly larger.
Thus time $5k+4$ has current $-t-2$ and full interval $[-t-2,2]$.
Distance $M$ now reaches $t+1$. Larger allowed positive jumps reach
at least $t+3$; the negative large jumps are farther away. The short
destinations $-t,-t+1$ are used, while $-t-4,-t-5$ have greater magnitude.
This proves the asserted next state.

Only the lane containing $t+1$ has its hull extended. If $t$ is odd,
its holes are $4,6,\ldots,t-1$; if $t$ is even, they are
$3,5,\ldots,t-1$. Their respective counts are $(t-3)/2$ and $(t-2)/2$,
giving the formula. Existence follows from unboundedness as before.
\end{proof}
These growing holes form one lane interval. The proposition therefore
does not obstruct a compact interval description, eventual coverage, or
a bound restricted to the prime distance set.

\begin{proposition}[Conditional single-lane filling]\label{prop:lane-fill}
Suppose 2 is the only positive even member of $D$. Let $B>L\ge1$ with
$B-L$ positive and even, and $\sigma\in\{-1,1\}$. If the current is
$\sigma B$ and the unused points inside $[-B,B]$ are exactly
$\{\sigma L,\sigma(L+2),\ldots,\sigma(B-2)\}$, then the next
$(B-L)/2$ moves are $\sigma(B-2),\sigma(B-4),\ldots,\sigma L$.
Any finite pattern of used points outside the interval is permitted.
\end{proposition}
\begin{proof}
All interior holes have the current parity. Odd distances cannot reach
them, and the only positive even distance is 2. Its inward destination
is the adjacent hole; its outward destination is used or exterior.
Every other unused reachable point is exterior and has magnitude greater
than $B$. The inward move is therefore the unique winner. Repeat until
no hole remains.
\end{proof}
The stated starting configuration and single-lane hypothesis do not follow merely
from allowing 2 and 3. None of these statements gives a bound on gaps between consecutive primes.

\section{Coverage through a moving boundary hole}\label{sec:repair}
Complete intervals are sufficient checkpoints, but are not necessary for coverage.
The following concrete rule has a persistent state shape with a changing hole.

\begin{theorem}[Coverage through a moving hole]\label{thm:repair}
Let $D=\{1,3,10\}\cup\{64+4j:j\ge0\}$. A unique infinite walk following the stated rule exists
and every such walk covers $\Z$ exactly once. At $n=72+4j$, put $b=36+2j$.
Then
\[
 a(n)=-b,\qquad U_n=([-b,b-1]\cap\Z)\cup\{b+1\},
 \qquad a(n+1)=b.
\]
At every time $n\ge64$ divisible by 4, every unused point of magnitude
at most $n/2$ is visited within three further moves. The bound is attained
at $n=64$; it is not asserted at arbitrary times or auxiliary states.
\end{theorem}
\begin{proof}
First prove a conditional transition. Suppose the state has current $-b$
and used set $([-b,b-1]\cap\Z)\cup\{b+1\}$, with $b\ge1$. Assume
$1,3,2b,2b+4$ are allowed and $2,2b+1,2b+2,2b+3$ are excluded.
The next four points are
\[
 b,\quad b+3,\quad -(b+1),\quad -(b+2).
\]
The first fills the least-magnitude hole by distance $2b$. The used set
then becomes $[-b,b+1]$. From $b$, the smaller unused competitors to
$b+3$ are $b+2,-(b+1),-(b+2)$, requiring the excluded distances
$2,2b+1,2b+2$. The negative tie loses, so distance 3 reaches $b+3$.
The allowed $2b+4$ crossing now reaches the least-magnitude unused point
$-(b+1)$. Finally the unit move reaches $-(b+2)$; its positive counterpart
would win the tie but requires the excluded distance $2b+3$.
The output is exactly the initial state shape with $b$ replaced by $b+2$.

Starting at $a(0)=0$, the first 72 destinations are the following list,
read from left to right and then from top to bottom:
\[
\begin{array}{rrrrrrrrrrrr}
1&2&-1&-2&-3&-4&-5&5&4&3&6&7\\
8&9&10&11&12&13&14&15&16&17&18&19\\
20&21&22&23&24&25&26&27&28&29&30&31\\
32&-32&-22&-12&-9&-6&-7&-8&-11&-10&-13&-14\\
-15&-16&-17&-18&-19&-20&-21&-24&-23&-26&-25&-28\\
-27&-30&-29&35&-33&-34&-31&33&34&37&-35&-36
\end{array}
\]
This is a finite certificate: for each proposed destination $v$, check
that it is unused, that $|v-c|\in D$ for the preceding current $c$, and
that every unused integer $w$ preceding $v$ in the priority order has
$|w-c|\notin D$. Only $|w|\le|v|$ need be checked. Direct application of
these tests verifies the displayed list; the same finite checks are proved
inside Lean in \code{RepairEntrance.lean}, supplied in the source package
\cite{package}. Uniqueness forces this prefix for every walk with this $D$.
In particular the list gives these exact used sets and currents:
\begin{align*}
 U_{64}&=([-32,32]\cap\Z)\setminus\{-31\}\ \cup\{35\},
   &a(64)&=35,\\
 U_{68}&=([-34,33]\cap\Z)\cup\{35\},&a(68)&=33,\\
 U_{72}&=([-36,35]\cap\Z)\cup\{37\},&a(72)&=-36.
\end{align*}
The next four destinations after time 64 are $-33,-34,-31,33$, and
$a(69)=34$. Thus the first two checkpoint holes are repaired at 67 and 69.

For $b=36+2j$, both $2b$ and $2b+4$ lie in the allowed tail. The three
intermediate distances are greater than 10 and not divisible by 4, so are
excluded, as is 2. The conditional transition therefore iterates from
time 72, proving the exact states and immediate repair thereafter.
Together with times 64 and 68 this proves the stated checkpoint bound.
The complete interior $[-36-2j,35+2j]$ expands without bound, establishing
coverage even though the symmetric checkpoint interval itself is never complete.
The rule forbids repeated visits, and unboundedness of $D$ gives existence.
\end{proof}
The local transition is reusable under its explicit premises, while the
coverage result for a walk starting at zero concerns one fixed initial path
and tail cutoff.
It does not prove the same initial state for other cutoffs or arbitrary sets
containing 1. The finite initial-path certificate and universal transition are
both mathematical proof obligations; an observed short repair delay alone
would establish neither.


\begin{thebibliography}{99}\small\raggedright
\bibitem{sloane2025} N. J. A. Sloane, \emph{In the Footsteps of Heraclitus},
DIMACS/Rutgers REU Seminar, 24 June 2025, p.~8.
\url{https://neilsloane.com/doc/REU2025.pdf}.
\bibitem{sloane2026} N. J. A. Sloane, \emph{Highlights from the First 400,000
Sequences in the OEIS}, Experimental Math Seminar, 10 September 2026, p.~9.
\url{https://neilsloane.com/doc/EMSep2026.pdf}.
\footnote{The date follows the PDF title page; the author's publication
list displays 2025 for this link.}
\bibitem{powers} R\'emy Sigrist, OEIS A379719,
powers-of-two Heraclitus walk and coverage description.
\url{https://oeis.org/A379719}.
\bibitem{fibonacci} G. Caveney, coverage announcement and six-block outline,
OEIS A377092, comment of 24 September 2026 referring to 21 September
SeqFan correspondence. \url{https://oeis.org/A377092}.
\bibitem{squares2} G. Caveney, OEIS A400196, squares-plus-2 coverage and
four-block outline, comment of 3 October 2026.
\url{https://oeis.org/A400196}.
\bibitem{adding2} G. Caveney, OEIS A400197, the role of adjoining distance 2.
\url{https://oeis.org/A400197}.
\bibitem{factorial} trureturing/The Omega Institute, \emph{Every integer
occurs in the factorial Heraclitus transform}, research dossier.
\url{https://www.trureturing.org/research/oeis-a393434-factorial-heraclitus/}.
\bibitem{factorialcode} trureturing/The Omega Institute,
\code{FactorialHeraclitusTransform.lean}, inspected repository commit
\code{535a6648df930c3bfdca8a945bfcebdb0d0e1ec7}, path
\code{D5/S3/Combinatorics/Permutation/}.
\url{https://github.com/the-omega-institute/trureturing}.
\bibitem{squares} N. J. A. Sloane and P. Xausa, proof announcement for the
square Heraclitus transform, OEIS A377091, 29 October 2026; under review.
\url{https://oeis.org/A377091}.
\bibitem{primes} R\'emy Sigrist and N. J. A. Sloane, OEIS A377090,
prime-distance Heraclitus walk and coverage conjecture.
\url{https://oeis.org/A377090}.
\bibitem{base7} R. Gujord, \emph{The base-seven Heraclitus transform:
manuscript and reproducible proof package}, version 1.0.0, Zenodo,
3 November 2026. \url{https://doi.org/10.5281/zenodo.23912682}.
\bibitem{package} R. Gujord, \emph{Coverage in Heraclitus walks: interval
filling and residue connectivity}, manuscript and reproducible proof package,
version 1.0.0, Zenodo, 2026.
\url{https://doi.org/10.5281/zenodo.23973601}.
\end{thebibliography}
\end{document}
